%=====================================================================
% PART I OPENER
%=====================================================================
\thispagestyle{plain}
% Part-opener maps get their own figure prefix (P1.1, ...).
\structuralfigures{P1.}

\begin{center}
\begin{tcolorbox}[enhanced, width=0.92\textwidth, colback=primary!5,
  colframe=primary, boxrule=0pt, leftrule=5pt, arc=3pt, top=8pt, bottom=8pt]
\large\itshape\color{primary!80!black}
Before any algorithm, three questions: what does it mean for a program to
\emph{learn}, what mathematics does learning run on, and what does a single
experiment look like from raw data to an honest number? Part I answers them, and
ends with the oldest precise answer to the first: learning as a search through
a space of hypotheses.
\end{tcolorbox}
\end{center}

\vspace{6pt}

\dsfigh{%
\begin{tikzpicture}[
  cbox/.style={rounded corners=3pt, draw=primary, line width=1pt, fill=primary!7,
               text=primary!80!black, font=\small\bfseries, align=center,
               minimum width=3.3cm, minimum height=1.0cm},
  hbox/.style={rounded corners=3pt, draw=primary, line width=1.8pt, fill=primary!16,
               text=primary!90!black, font=\small\bfseries, align=center,
               minimum width=3.3cm, minimum height=1.0cm},
  arr/.style={-{Stealth[length=2.4mm]}, primary!60, line width=1pt}
]
  \node[cbox] (c1) at (0,0)    {Ch.\ 1\\Introduction};
  \node[cbox] (c2) at (4.2,0)  {Ch.\ 2\\Mathematics};
  \node[hbox] (c3) at (8.4,0)  {Ch.\ 3\\Workflow \& Toolkit};
  \node[cbox] (c4) at (12.6,0) {Ch.\ 4\\Concept Learning};
  \draw[arr] (c1) -- (c2);
  \draw[arr] (c2) -- (c3);
  \draw[arr] (c3) -- (c4);
  \draw[arr] (c1.south) .. controls +(0,-1.3) and +(0,-1.3) .. (c4.south);
  \node[font=\scriptsize\itshape, text=primary, align=center, fill=white,
        inner sep=2pt] at (6.3,-1.0)
       {the hypothesis-space view of Chapter 1, made precise};
  \node[font=\scriptsize\itshape, text=accent, align=center] at (8.4,0.85)
       {every lab in the course uses it};
\end{tikzpicture}}%
{Reading path through Part I. Chapter 3 (highlighted) sets up the workflow that every
later lab repeats; Chapter 4 formalises the idea of learning introduced in Chapter 1.}{part1map}

\newpage
